\documentclass[12pt]{article}
\usepackage{times}
\usepackage[margin=1in]{geometry}
\usepackage{graphicx}
\usepackage{amsmath}
\usepackage[justification=raggedright,singlelinecheck=false,font=small]{caption}
\usepackage[colorlinks=true,linkcolor=blue,citecolor=blue,urlcolor=blue]{hyperref}
\usepackage{setspace}
\usepackage{booktabs}

\setstretch{1.15}
\sloppy

\title{Differential Gene Expression Analysis of Breast Tumor and Normal Tissue \\
Using Bulk RNA-Sequencing Data}
\date{}
\begin{document}
\maketitle

\begin{abstract}
\noindent RNA sequencing (RNA-seq) enables high-throughput quantification of gene
expression across the transcriptome. Here, differential expression analysis was
performed on a public breast cancer RNA-seq dataset (GEO accession GSE183947,
$n=60$; 30 tumor and 30 matched normal samples) to identify differentially
expressed genes (DEGs) between tumor and normal tissue. Significant DEGs were
identified using DESeq2, characterized functionally via Gene Ontology (GO) and
KEGG pathway enrichment, and evaluated as classification features using a
Random Forest model.
\end{abstract}

\section{Introduction}
RNA-seq has revolutionized transcriptome analysis by enabling high-throughput quantification of gene expression across the entire genome (1). Compared to microarrays, RNA-seq offers greater sensitivity and the ability to detect novel transcripts and alternative splicing events (2). These advantages make RNA-seq an indispensable tool for investigating the molecular mechanisms underlying health and disease.
A central application of RNA-seq is the identification of differentially expressed genes (DEGs), which reflect transcriptional alterations between biological states such as tumor and normal tissue (3).

\subsection{FPKM normalization}
FPKM (Fragments Per Kilobase of transcript per Million mapped reads) is a
normalization metric commonly used in RNA-seq experiments. The number of reads
mapped to a gene reflects its expression level, but raw read counts cannot be
directly compared across genes or samples due to differences in sequencing
depth and gene length. FPKM normalizes expression data by accounting for both:
\begin{itemize}
  \item \textbf{Sequencing depth} - dividing by the total number of mapped reads.
  \item \textbf{Gene length} - dividing by transcript length in kilobases.
\end{itemize}

\section{Dataset}
The RNA-seq dataset GSE183947 was retrieved from the NCBI Gene Expression
Omnibus (GEO) database. This dataset comprises 60 samples of breast tissue,
including 30 breast tumor and 30 normal tissue samples. The dataset was
already pre-normalized in FPKM units, which were rounded and transformed into
integer counts before DESeq2 analysis.

\section{Tools Required}
\begin{itemize}
  \item \textbf{R programming} - used to perform all steps of differential
        expression analysis (DEG identification, normalization, statistical
        testing), functional enrichment, and visualization.
  \item \textbf{DESeq2} - for identifying DEGs from count-based sequence data.
  \item \textbf{clusterProfiler} - functional enrichment analysis, including
        Gene Ontology (GO) and KEGG pathway analysis.
  \item \textbf{enrichR} - disease association, tissue enrichment, and
        pathway analysis.
  \item \textbf{ggplot2} - boxplots, volcano plots, MA plots.
  \item \textbf{pheatmap} - clustered heatmaps of DEGs.
  \item \textbf{randomForest} - machine learning models using normalized
        expression values of top DEGs to classify/rank important genes.
  \item \textbf{GEOquery} - downloading the expression matrix and metadata
        from the GEO database.
\end{itemize}

\section{Methodology}

\subsection{Data acquisition and preparation}
The RNA-seq dataset (60 samples: 30 tumor, 30 normal) was imported into R.
Phenotypic and experimental annotations for the dataset were retrieved using
the \texttt{pData(phenoData(...))} function. Metadata provided group
information differentiating 30 tumor versus 30 normal samples.

\subsection{Metadata processing and reshaping}
Relevant columns were selected and renamed for clarity (tissue and metastasis
status). The raw FPKM matrix was converted from wide format to long format
using \texttt{gather()}, creating a table with one row per gene-sample pair.
The long-format expression data was then merged with the cleaned metadata
based on sample identifiers.

\subsection{Count matrix and sample information}
DESeq2 requires integer counts. FPKM values were rounded to the nearest
integer to generate an approximate count matrix. The cleaned metadata was used
to create a sample information table, where each sample was assigned to a
condition label: Tumor or Normal. Sample ordering between the count matrix and
sample information table was explicitly verified before model construction.

\subsection{Differential gene expression analysis}
A \texttt{DESeqDataSet} was constructed from the integer count matrix and
sample information table, with the design formula $\sim \text{condition}$.
Genes with extremely low expression across samples (row sum $< 10$) were
filtered out to reduce noise. DESeq2's internal normalization procedure was
applied, and a negative binomial generalized linear model (GLM) was fit to
each gene. Genes with an adjusted $p$-value (Benjamini--Hochberg) $\leq 0.05$
and an absolute $\log_2$ fold change $\geq 1$ were considered significantly
differentially expressed.

\subsection{Visualization}
Boxplots of the top five DEGs (ranked by adjusted $p$-value) were generated
using DESeq2-normalized, $\log_2$-transformed counts, stratified by condition.
A volcano plot was constructed by plotting $\log_2$ fold change against
$-\log_{10}$(adjusted $p$-value), highlighting upregulated and downregulated
genes. An MA plot and a clustered heatmap of the top 50 DEGs (Euclidean
distance, complete linkage, row-scaled) were additionally generated to assess
overall expression patterns between conditions.

\subsection{Functional enrichment analysis}
Gene Ontology (GO) Biological Process enrichment was performed on upregulated
genes using \texttt{clusterProfiler}. KEGG pathway and disease/tissue
enrichment analyses were additionally performed using \texttt{enrichR} on both
upregulated and downregulated gene sets.

\subsection{Machine learning-based classification using top DEGs}
\label{sec:ml}
To evaluate the discriminative potential of the identified DEGs, a Random
Forest classifier was trained using a nested design in which DEG selection
and test-set evaluation are kept fully separated, ensuring that no
information from held-out samples influences which genes are chosen or how
expression values are normalized.

\begin{enumerate}
  \item \textbf{Train/test split:} All 60 samples were partitioned into 70\%
        training and 30\% testing sets using stratified sampling
        (\texttt{caret::createDataPartition}), performed \emph{before} any
        differential expression analysis for the classification pipeline.
  \item \textbf{Training-only DEG selection:} DESeq2 was run using only the
        training samples; the top 10 genes (ranked by adjusted $p$-value)
        from this training-only fit were used as classifier features. This
        keeps feature selection blind to the test partition, so that the
        genes chosen are not influenced by the very samples later used to
        evaluate the model.
  \item \textbf{Feature construction:} Training-set features were the
        DESeq2-normalized, $\log_2$-transformed counts from the
        training-only fit, joined to samples by sample identifier (rather
        than row position, which would risk silent sample/label
        misalignment). Test-set samples were normalized independently,
        against the training set's reference geometric means
        (via \texttt{estimateSizeFactorsForMatrix()}), so that no
        information from the test set - including its own normalization - entered feature selection or scaling at any point.
  \item \textbf{Model training:} A Random Forest classifier was trained with
        500 trees ($n_{\text{tree}} = 500$) on the training features.
  \item \textbf{Model evaluation:} Predicted class labels were generated for
        the held-out test set, and classification performance was assessed
        using a confusion matrix.
  \item \textbf{Feature importance:} Gene importance was computed using the
        Mean Decrease in Gini index from the trained model.
  \item \textbf{Visualization:} Feature importance was visualized as a
        horizontal bar plot using \texttt{ggplot2}.
  \item \textbf{Cross-validation robustness check:} Because a single 70/30
        split is evaluated on only $\sim$18 test samples - where a single misclassification changes accuracy by roughly six percentage points - the identical nested procedure (stratified split, training-only DEG selection, training-referenced test normalization, Random Forest training and evaluation) was additionally repeated across
        5-fold cross-validation with 5 repeats (25 total fold fits), independently reselecting the top DEGs within each fold's training partition. Mean accuracy, sensitivity, and specificity across folds, together with their standard deviations, were reported alongside the single-split result to characterize the variability of the
        performance estimate.
\end{enumerate}

Note that the DESeq2 results reported in the differential expression and
enrichment sections above (Sections~4--6) use the full 60-sample dataset, as
is standard practice for exploratory DEG discovery; the nested, sample-split
design described here applies specifically to the machine learning
evaluation, where a genuinely held-out accuracy estimate is required.

\section*{References}
\begin{enumerate}
  \item Wang, Z., Gerstein, M., Snyder, M. RNA-Seq: a revolutionary tool for
  transcriptomics. Nature Reviews Genetics. 2009;10:57--63.
  \item Ozsolak, F., Milos, P.M. RNA sequencing: advances, challenges and
  opportunities. Nature Reviews Genetics. 2011;12:87--98.
  \item Love, M.I., Huber, W., Anders, S. Moderated estimation of fold change
  and dispersion for RNA-seq data with DESeq2. Genome Biology.
  2014;15:550.
\end{enumerate}

\end{document}
